\section{Conclusion}
\label{sec:conclusion}

We have shown that within-game z-score normalization is the correct preprocessing step
for identifying independent dimensions of game design character. Applied to the 150-game
TIGRIS corpus, this normalization reduces PC1 from 81\% (raw) to $\leq 55\%$ and reveals
4--5 genuinely independent components explaining $\geq 80\%$ of within-game normalized
variance. The five components are labeled Tension, Interaction, Game-Architecture,
Elegance, and Range --- the TIGER dimensions formalized in TG-B1.

The methodological lesson is general: any domain where quality is a dominant latent
factor in multi-attribute scores will exhibit a strong general PC that obscures
within-domain variation. Within-group normalization before PCA is the standard remedy,
but it has not previously been applied to game design scoring. We recommend this approach
for future multi-axis evaluation studies in creative domains.

All analysis code and PCA outputs are available at [URL].
